%%%%%%%%%%%%%%%%%%%%%%%%%%%%%%%%%%%%%%%%%%%%%%%%%%%%%%%%%%%%%%%%%%%%%%%%%%%%%%%%
% CHAPTER 1
%%%%%%%%%%%%%%%%%%%%%%%%%%%%%%%%%%%%%%%%%%%%%%%%%%%%%%%%%%%%%%%%%%%%%%%%%%%%%%%%

\chapter{Introduction and Classification of Signals}

\begin{tcolorbox}[colback=blue!5!white,colframe=blue!70!black,title=Learning Objectives]

After studying this chapter, the reader will be able to

\begin{itemize}
\item Define a signal and understand its role in engineering.
\item Distinguish between continuous-time and discrete-time signals.
\item Classify signals as analog or digital.
\item Identify deterministic and random signals.
\item Understand periodic and aperiodic signals.
\item Recognize energy and power signals.
\item Determine whether a signal is even or odd.
\item Understand elementary signals used in Signals and Systems.
\end{itemize}

\end{tcolorbox}

%%%%%%%%%%%%%%%%%%%%%%%%%%%%%%%%%%%%%%%%%%%%%%%%%%%%%%%%%%%%%%%%%%%%%%%%%%%%%%%%

\section{Introduction}

Signals and Systems is one of the most fundamental subjects in Electronics and
Communication Engineering. Every communication system, control system, embedded
system and signal-processing algorithm is built upon the concepts of signals and
their transformation through systems.

Examples include

\begin{itemize}
\item speech signals,
\item ECG signals,
\item satellite telemetry,
\item sensor outputs,
\item image signals,
\item and digital communication waveforms.
\end{itemize}

The main objective of this subject is to understand

\begin{enumerate}
\item how signals are represented,
\item how systems operate on signals,
\item and how signals can be analyzed in time and frequency domains.
\end{enumerate}

%%%%%%%%%%%%%%%%%%%%%%%%%%%%%%%%%%%%%%%%%%%%%%%%%%%%%%%%%%%%%%%%%%%%%%%%%%%%%%%%

\section{What is a Signal?}

\subsection{Definition}

A \textbf{signal} is a function that conveys information about the behavior or
nature of a phenomenon.

Mathematically,

\[
x(t),\qquad x[n],\qquad x(t_1,t_2,\ldots)
\]

are all signals.

Here,

\begin{itemize}
\item \(x\) is the dependent variable (signal value),
\item \(t\) or \(n\) is the independent variable.
\end{itemize}

%%%%%%%%%%%%%%%%%%%%%%%%%%%%%%%%%%%%%%%%%%%%%%%%%%%%%%%%%%%%%%%%%%%%%%%%%%%%%%%%

\subsection{Examples}

\begin{table}[H]
\centering
\caption{Examples of engineering signals}
\begin{tabular}{|p{5cm}|p{6cm}|}
\hline
\textbf{Signal} & \textbf{Mathematical Form} \\
\hline
DC voltage & \(x(t)=5\) \\
\hline
Sinusoidal voltage & \(x(t)=10\sin(100\pi t)\) \\
\hline
Temperature variation & \(x(t)=25+2\sin(0.1t)\) \\
\hline
Digital bit sequence & \(x[n]\in\{0,1\}\) \\
\hline
\end{tabular}
\end{table}

%%%%%%%%%%%%%%%%%%%%%%%%%%%%%%%%%%%%%%%%%%%%%%%%%%%%%%%%%%%%%%%%%%%%%%%%%%%%%%%%

\section{Independent Variable}

The independent variable may represent

\begin{itemize}
\item time,
\item position,
\item angle,
\item frequency,
\item or any measurable quantity.
\end{itemize}

In Signals and Systems, the most common independent variable is \textbf{time}.

%%%%%%%%%%%%%%%%%%%%%%%%%%%%%%%%%%%%%%%%%%%%%%%%%%%%%%%%%%%%%%%%%%%%%%%%%%%%%%%%

\section{Continuous-Time Signal}

\subsection{Definition}

A signal defined for every value of time is called a
\textbf{continuous-time (CT) signal}.

\[
x(t),\qquad -\infty<t<\infty
\]

%%%%%%%%%%%%%%%%%%%%%%%%%%%%%%%%%%%%%%%%%%%%%%%%%%%%%%%%%%%%%%%%%%%%%%%%%%%%%%%%

\subsection{Example}

\[
x(t)=A\cos(\omega t+\phi)
\]

%%%%%%%%%%%%%%%%%%%%%%%%%%%%%%%%%%%%%%%%%%%%%%%%%%%%%%%%%%%%%%%%%%%%%%%%%%%%%%%%

\subsection{Graph}

\begin{center}
\begin{tikzpicture}[scale=0.9]
\draw[->] (-0.5,0) -- (6.5,0) node[right] {$t$};
\draw[->] (0,-1.5) -- (0,1.5) node[above] {$x(t)$};

\draw[thick,smooth,samples=100,domain=0:6]
plot(\x,{sin(deg(1.2*\x))});

\node at (3,1.2) {$x(t)$};
\end{tikzpicture}
\end{center}

The signal exists for all time instants.

%%%%%%%%%%%%%%%%%%%%%%%%%%%%%%%%%%%%%%%%%%%%%%%%%%%%%%%%%%%%%%%%%%%%%%%%%%%%%%%%

\section{Discrete-Time Signal}

\subsection{Definition}

A signal defined only at discrete instants of time is called a
\textbf{discrete-time (DT) signal}.

\[
x[n],\qquad n\in\mathbb{Z}
\]

%%%%%%%%%%%%%%%%%%%%%%%%%%%%%%%%%%%%%%%%%%%%%%%%%%%%%%%%%%%%%%%%%%%%%%%%%%%%%%%%

\subsection{Example}

\[
x[n]=\sin\left(\frac{\pi n}{4}\right)
\]

%%%%%%%%%%%%%%%%%%%%%%%%%%%%%%%%%%%%%%%%%%%%%%%%%%%%%%%%%%%%%%%%%%%%%%%%%%%%%%%%

\subsection{Graph}

\begin{center}
\begin{tikzpicture}[scale=0.9]
\draw[->] (-0.5,0) -- (7.5,0) node[right] {$n$};
\draw[->] (0,-1.5) -- (0,1.5) node[above] {$x[n]$};

\foreach \n in {0,1,2,3,4,5,6,7}{
  \pgfmathsetmacro{\y}{sin(deg(pi*\n/4))}
  \draw[thick] (\n,0) -- (\n,\y);
  \filldraw (\n,\y) circle (2pt);
}
\end{tikzpicture}
\end{center}

Only the sampled values exist.

%%%%%%%%%%%%%%%%%%%%%%%%%%%%%%%%%%%%%%%%%%%%%%%%%%%%%%%%%%%%%%%%%%%%%%%%%%%%%%%%

\section{Analog and Digital Signals}

%%%%%%%%%%%%%%%%%%%%%%%%%%%%%%%%%%%%%%%%%%%%%%%%%%%%%%%%%%%%%%%%%%%%%%%%%%%%%%%%

\subsection{Analog Signal}

An analog signal has a continuous range of amplitudes.

\[
x(t)\in\mathbb{R}
\]

Examples:

\begin{itemize}
\item microphone output,
\item ECG waveform,
\item temperature sensor voltage.
\end{itemize}

%%%%%%%%%%%%%%%%%%%%%%%%%%%%%%%%%%%%%%%%%%%%%%%%%%%%%%%%%%%%%%%%%%%%%%%%%%%%%%%%

\subsection{Digital Signal}

A digital signal has a finite set of amplitude values.

For binary signals,

\[
x[n]\in\{0,1\}
\]

Examples:

\begin{itemize}
\item UART data,
\item SPI signals,
\item microcontroller GPIO levels.
\end{itemize}

%%%%%%%%%%%%%%%%%%%%%%%%%%%%%%%%%%%%%%%%%%%%%%%%%%%%%%%%%%%%%%%%%%%%%%%%%%%%%%%%

\section{Deterministic and Random Signals}

%%%%%%%%%%%%%%%%%%%%%%%%%%%%%%%%%%%%%%%%%%%%%%%%%%%%%%%%%%%%%%%%%%%%%%%%%%%%%%%%

\subsection{Deterministic Signal}

A signal whose value can be predicted exactly.

Example:

\[
x(t)=5\cos(2\pi 100t)
\]

%%%%%%%%%%%%%%%%%%%%%%%%%%%%%%%%%%%%%%%%%%%%%%%%%%%%%%%%%%%%%%%%%%%%%%%%%%%%%%%%

\subsection{Random Signal}

A signal whose future values cannot be predicted exactly.

Examples:

\begin{itemize}
\item thermal noise,
\item atmospheric noise,
\item stock-price fluctuations.
\end{itemize}

%%%%%%%%%%%%%%%%%%%%%%%%%%%%%%%%%%%%%%%%%%%%%%%%%%%%%%%%%%%%%%%%%%%%%%%%%%%%%%%%

\section{Periodic Signal}

\subsection{Continuous-Time}

A signal is periodic if

\[
\boxed{
x(t)=x(t+T)
}
\]

for some positive constant \(T\).

The smallest such \(T\) is called the
\textbf{fundamental period}.

Example:

\[
x(t)=\sin(2\pi f t)
\]

\[
T=\frac{1}{f}
\]

%%%%%%%%%%%%%%%%%%%%%%%%%%%%%%%%%%%%%%%%%%%%%%%%%%%%%%%%%%%%%%%%%%%%%%%%%%%%%%%%

\subsection{Discrete-Time}

A discrete-time signal is periodic if

\[
\boxed{
x[n]=x[n+N]
}
\]

where \(N\) is the fundamental period in samples.

Example:

\[
x[n]=\cos\left(\frac{\pi n}{2}\right)
\]

Here,

\[
N=4.
\]

%%%%%%%%%%%%%%%%%%%%%%%%%%%%%%%%%%%%%%%%%%%%%%%%%%%%%%%%%%%%%%%%%%%%%%%%%%%%%%%%

\section{Aperiodic Signal}

A signal that does not repeat is called \textbf{aperiodic}.

Example:

\[
x(t)=e^{-t}u(t)
\]

This signal decays with time and never repeats.

%%%%%%%%%%%%%%%%%%%%%%%%%%%%%%%%%%%%%%%%%%%%%%%%%%%%%%%%%%%%%%%%%%%%%%%%%%%%%%%%

\section{Even and Odd Signals}

%%%%%%%%%%%%%%%%%%%%%%%%%%%%%%%%%%%%%%%%%%%%%%%%%%%%%%%%%%%%%%%%%%%%%%%%%%%%%%%%

\subsection{Even Signal}

\[
\boxed{
x(t)=x(-t)
}
\]

Example:

\[
x(t)=t^2
\]

%%%%%%%%%%%%%%%%%%%%%%%%%%%%%%%%%%%%%%%%%%%%%%%%%%%%%%%%%%%%%%%%%%%%%%%%%%%%%%%%

\subsection{Odd Signal}

\[
\boxed{
x(t)=-x(-t)
}
\]

Example:

\[
x(t)=t
\]

%%%%%%%%%%%%%%%%%%%%%%%%%%%%%%%%%%%%%%%%%%%%%%%%%%%%%%%%%%%%%%%%%%%%%%%%%%%%%%%%

\section{Signal Decomposition}

Any signal can be written as

\[
\boxed{
x(t)=x_e(t)+x_o(t)
}
\]

where

\[
x_e(t)=\frac{x(t)+x(-t)}{2}
\]

\[
x_o(t)=\frac{x(t)-x(-t)}{2}
\]

This decomposition is very important in Fourier analysis.

%%%%%%%%%%%%%%%%%%%%%%%%%%%%%%%%%%%%%%%%%%%%%%%%%%%%%%%%%%%%%%%%%%%%%%%%%%%%%%%%

\begin{tcolorbox}[title=Quick Recall,colback=yellow!10]

\[
x(t)\rightarrow \text{continuous-time}
\]

\[
x[n]\rightarrow \text{discrete-time}
\]

\[
x(t)=x(t+T)\rightarrow \text{periodic}
\]

\[
x(t)=x(-t)\rightarrow \text{even}
\]

\[
x(t)=-x(-t)\rightarrow \text{odd}
\]

\end{tcolorbox}
%%%%%%%%%%%%%%%%%%%%%%%%%%%%%%%%%%%%%%%%%%%%%%%%%%%%%%%%%%%%%%%%%%%%%%%%%%%%%%%%
%%%%%%%%%%%%%%%%%%%%%%%%%%%%%%%%%%%%%%%%%%%%%%%%%%%%%%%%%%%%%%%%%%%%%%%%%%%%%%%%
\section{Energy and Power Signals}

%%%%%%%%%%%%%%%%%%%%%%%%%%%%%%%%%%%%%%%%%%%%%%%%%%%%%%%%%%%%%%%%%%%%%%%%%%%%%%%%

\subsection{Energy of a Continuous-Time Signal}

The energy of a signal is

\[
\boxed{
E=\int_{-\infty}^{\infty}|x(t)|^2dt
}
\]

A signal is called an \textbf{energy signal} if

\[
0<E<\infty.
\]

%%%%%%%%%%%%%%%%%%%%%%%%%%%%%%%%%%%%%%%%%%%%%%%%%%%%%%%%%%%%%%%%%%%%%%%%%%%%%%%%

\subsection{Example}

Consider

\[
x(t)=e^{-t}u(t).
\]

Then,

\[
E=\int_0^\infty e^{-2t}dt
\]

\[
E=\left[-\frac{e^{-2t}}{2}\right]_0^\infty
\]

\[
\boxed{
E=\frac{1}{2}
}
\]

Hence, this is an energy signal.

%%%%%%%%%%%%%%%%%%%%%%%%%%%%%%%%%%%%%%%%%%%%%%%%%%%%%%%%%%%%%%%%%%%%%%%%%%%%%%%%

\section{Average Power of a Continuous-Time Signal}

\[
\boxed{
P=\lim_{T\to\infty}\frac{1}{2T}
\int_{-T}^{T}|x(t)|^2dt
}
\]

A signal is called a \textbf{power signal} if

\[
0<P<\infty.
\]

%%%%%%%%%%%%%%%%%%%%%%%%%%%%%%%%%%%%%%%%%%%%%%%%%%%%%%%%%%%%%%%%%%%%%%%%%%%%%%%%

\subsection{Example}

For

\[
x(t)=A\cos(\omega t),
\]

\[
P=\frac{A^2}{2}.
\]

Thus, a sinusoidal signal is a power signal.

%%%%%%%%%%%%%%%%%%%%%%%%%%%%%%%%%%%%%%%%%%%%%%%%%%%%%%%%%%%%%%%%%%%%%%%%%%%%%%%%

\section{Relationship Between Energy and Power}

\begin{table}[H]
\centering
\caption{Energy and power signals}
\begin{tabular}{|p{5cm}|p{5cm}|}
\hline
\textbf{Energy Signal} & \textbf{Power Signal} \\
\hline
Finite energy & Infinite energy \\
\hline
Zero average power & Finite average power \\
\hline
Usually non-periodic & Usually periodic \\
\hline
Example: pulse & Example: sinusoid \\
\hline
\end{tabular}
\end{table}

A non-zero signal cannot be both an energy signal and a power signal.

%%%%%%%%%%%%%%%%%%%%%%%%%%%%%%%%%%%%%%%%%%%%%%%%%%%%%%%%%%%%%%%%%%%%%%%%%%%%%%%%

\section{Unit Impulse Signal}

%%%%%%%%%%%%%%%%%%%%%%%%%%%%%%%%%%%%%%%%%%%%%%%%%%%%%%%%%%%%%%%%%%%%%%%%%%%%%%%%

\subsection{Definition}

The \textbf{unit impulse} is denoted by

\[
\boxed{
\delta(t)
}
\]

and is defined by

\[
\delta(t)=0,\qquad t\neq0
\]

with the property

\[
\boxed{
\int_{-\infty}^{\infty}\delta(t)\,dt=1
}
\]

%%%%%%%%%%%%%%%%%%%%%%%%%%%%%%%%%%%%%%%%%%%%%%%%%%%%%%%%%%%%%%%%%%%%%%%%%%%%%%%%

\subsection{Graph}

\begin{center}
\begin{tikzpicture}[scale=0.9]
\draw[->] (-2,0) -- (2,0) node[right] {$t$};
\draw[->] (0,0) -- (0,2) node[above] {$\delta(t)$};

\draw[very thick] (0,0) -- (0,1.8);
\filldraw (0,1.8) circle (2pt);
\node[right] at (0,1.6) {Area = 1};
\end{tikzpicture}
\end{center}

%%%%%%%%%%%%%%%%%%%%%%%%%%%%%%%%%%%%%%%%%%%%%%%%%%%%%%%%%%%%%%%%%%%%%%%%%%%%%%%%

\subsection{Sifting Property}

\[
\boxed{
\int_{-\infty}^{\infty}x(t)\delta(t-t_0)\,dt=x(t_0)
}
\]

The impulse extracts the value of the signal at \(t=t_0\).

%%%%%%%%%%%%%%%%%%%%%%%%%%%%%%%%%%%%%%%%%%%%%%%%%%%%%%%%%%%%%%%%%%%%%%%%%%%%%%%%

\section{Unit Step Signal}

%%%%%%%%%%%%%%%%%%%%%%%%%%%%%%%%%%%%%%%%%%%%%%%%%%%%%%%%%%%%%%%%%%%%%%%%%%%%%%%%

\subsection{Definition}

The \textbf{unit step} is denoted by

\[
\boxed{
u(t)
}
\]

and is defined as

\[
u(t)=
\begin{cases}
0, & t<0\\
1, & t>0
\end{cases}
\]

%%%%%%%%%%%%%%%%%%%%%%%%%%%%%%%%%%%%%%%%%%%%%%%%%%%%%%%%%%%%%%%%%%%%%%%%%%%%%%%%

\subsection{Graph}

\begin{center}
\begin{tikzpicture}[scale=0.9]
\draw[->] (-3,0) -- (3,0) node[right] {$t$};
\draw[->] (0,-0.5) -- (0,1.8) node[above] {$u(t)$};

\draw[thick] (-3,0) -- (0,0);
\draw[thick] (0,1) -- (3,1);
\filldraw (0,1) circle (2pt);
\end{tikzpicture}
\end{center}

%%%%%%%%%%%%%%%%%%%%%%%%%%%%%%%%%%%%%%%%%%%%%%%%%%%%%%%%%%%%%%%%%%%%%%%%%%%%%%%%

\subsection{Relationship with Impulse}

\[
\boxed{
\frac{d}{dt}u(t)=\delta(t)
}
\]

and

\[
\boxed{
u(t)=\int_{-\infty}^{t}\delta(\tau)\,d\tau
}
\]

%%%%%%%%%%%%%%%%%%%%%%%%%%%%%%%%%%%%%%%%%%%%%%%%%%%%%%%%%%%%%%%%%%%%%%%%%%%%%%%%

\section{Delayed Step Signal}

A delayed step is

\[
\boxed{
u(t-t_0)
}
\]

It becomes 1 at \(t=t_0\).

%%%%%%%%%%%%%%%%%%%%%%%%%%%%%%%%%%%%%%%%%%%%%%%%%%%%%%%%%%%%%%%%%%%%%%%%%%%%%%%%

\section{Rectangular Pulse}

A pulse of width \(T\) is represented as

\[
\boxed{
p(t)=u(t)-u(t-T)
}
\]

%%%%%%%%%%%%%%%%%%%%%%%%%%%%%%%%%%%%%%%%%%%%%%%%%%%%%%%%%%%%%%%%%%%%%%%%%%%%%%%%

\subsection{Graph}

\begin{center}
\begin{tikzpicture}[scale=0.9]
\draw[->] (-1,0) -- (6,0) node[right] {$t$};
\draw[->] (0,-0.5) -- (0,1.8) node[above] {$p(t)$};

\draw[thick] (-1,0) -- (0,0);
\draw[thick] (0,1) -- (4,1);
\draw[thick] (4,1) -- (4,0);
\draw[thick] (4,0) -- (6,0);

\node[below] at (0,0) {$0$};
\node[below] at (4,0) {$T$};
\end{tikzpicture}
\end{center}

%%%%%%%%%%%%%%%%%%%%%%%%%%%%%%%%%%%%%%%%%%%%%%%%%%%%%%%%%%%%%%%%%%%%%%%%%%%%%%%%

\section{Ramp Signal}

%%%%%%%%%%%%%%%%%%%%%%%%%%%%%%%%%%%%%%%%%%%%%%%%%%%%%%%%%%%%%%%%%%%%%%%%%%%%%%%%

\subsection{Definition}

The ramp signal is

\[
\boxed{
r(t)=tu(t)
}
\]

%%%%%%%%%%%%%%%%%%%%%%%%%%%%%%%%%%%%%%%%%%%%%%%%%%%%%%%%%%%%%%%%%%%%%%%%%%%%%%%%

\subsection{Graph}

\begin{center}
\begin{tikzpicture}[scale=0.9]
\draw[->] (-1,0) -- (5,0) node[right] {$t$};
\draw[->] (0,-0.5) -- (0,4) node[above] {$r(t)$};

\draw[thick] (-1,0) -- (0,0);
\draw[thick] (0,0) -- (4,3.2);
\end{tikzpicture}
\end{center}

%%%%%%%%%%%%%%%%%%%%%%%%%%%%%%%%%%%%%%%%%%%%%%%%%%%%%%%%%%%%%%%%%%%%%%%%%%%%%%%%

\subsection{Derivative}

\[
\boxed{
\frac{d}{dt}r(t)=u(t)
}
\]

Thus,

\[
r(t)\xrightarrow{d/dt}u(t)\xrightarrow{d/dt}\delta(t).
\]

%%%%%%%%%%%%%%%%%%%%%%%%%%%%%%%%%%%%%%%%%%%%%%%%%%%%%%%%%%%%%%%%%%%%%%%%%%%%%%%%

\section{Exponential Signal}

%%%%%%%%%%%%%%%%%%%%%%%%%%%%%%%%%%%%%%%%%%%%%%%%%%%%%%%%%%%%%%%%%%%%%%%%%%%%%%%%

\subsection{General Form}

\[
\boxed{
x(t)=Ae^{at}
}
\]

%%%%%%%%%%%%%%%%%%%%%%%%%%%%%%%%%%%%%%%%%%%%%%%%%%%%%%%%%%%%%%%%%%%%%%%%%%%%%%%%

\subsection{Cases}

\begin{table}[H]
\centering
\caption{Exponential signal behavior}
\begin{tabular}{|c|c|}
\hline
\(a\) & Behavior \\
\hline
\(a<0\) & Decaying exponential \\
\hline
\(a>0\) & Growing exponential \\
\hline
\(a=0\) & Constant signal \\
\hline
\end{tabular}
\end{table}

%%%%%%%%%%%%%%%%%%%%%%%%%%%%%%%%%%%%%%%%%%%%%%%%%%%%%%%%%%%%%%%%%%%%%%%%%%%%%%%%

\subsection{Decaying Exponential}

\[
x(t)=e^{-at}u(t),\qquad a>0
\]

%%%%%%%%%%%%%%%%%%%%%%%%%%%%%%%%%%%%%%%%%%%%%%%%%%%%%%%%%%%%%%%%%%%%%%%%%%%%%%%%

\subsection{Graph}

\begin{center}
\begin{tikzpicture}[scale=0.9]
\draw[->] (-0.5,0) -- (6,0) node[right] {$t$};
\draw[->] (0,0) -- (0,3) node[above] {$x(t)$};

\draw[thick,smooth,samples=100,domain=0:5.5]
plot(\x,{2.5*exp(-0.6*\x)});
\end{tikzpicture}
\end{center}

%%%%%%%%%%%%%%%%%%%%%%%%%%%%%%%%%%%%%%%%%%%%%%%%%%%%%%%%%%%%%%%%%%%%%%%%%%%%%%%%

\section{Sinusoidal Signal}

%%%%%%%%%%%%%%%%%%%%%%%%%%%%%%%%%%%%%%%%%%%%%%%%%%%%%%%%%%%%%%%%%%%%%%%%%%%%%%%%

\subsection{General Form}

\[
\boxed{
x(t)=A\cos(\omega t+\phi)
}
\]

where

\begin{itemize}
\item \(A\): amplitude,
\item \(\omega\): angular frequency (rad/s),
\item \(\phi\): phase angle.
\end{itemize}

%%%%%%%%%%%%%%%%%%%%%%%%%%%%%%%%%%%%%%%%%%%%%%%%%%%%%%%%%%%%%%%%%%%%%%%%%%%%%%%%

\subsection{Frequency and Period}

\[
\boxed{
f=\frac{\omega}{2\pi}
}
\]

\[
\boxed{
T=\frac{1}{f}=\frac{2\pi}{\omega}
}
\]

%%%%%%%%%%%%%%%%%%%%%%%%%%%%%%%%%%%%%%%%%%%%%%%%%%%%%%%%%%%%%%%%%%%%%%%%%%%%%%%%

\subsection{Graph}

\begin{center}
\begin{tikzpicture}[scale=0.9]
\draw[->] (-0.5,0) -- (7,0) node[right] {$t$};
\draw[->] (0,-2) -- (0,2) node[above] {$x(t)$};

\draw[thick,smooth,samples=100,domain=0:6.5]
plot(\x,{1.5*cos(deg(1.2*\x))});
\end{tikzpicture}
\end{center}

%%%%%%%%%%%%%%%%%%%%%%%%%%%%%%%%%%%%%%%%%%%%%%%%%%%%%%%%%%%%%%%%%%%%%%%%%%%%%%%%

\section{Complex Exponential Signal}

Using Euler's formula,

\[
\boxed{
e^{j\theta}=\cos\theta+j\sin\theta
}
\]

Therefore,

\[
\boxed{
x(t)=Ae^{j\omega t}
}
\]

is a complex sinusoid.

This representation is extensively used in Fourier analysis.

%%%%%%%%%%%%%%%%%%%%%%%%%%%%%%%%%%%%%%%%%%%%%%%%%%%%%%%%%%%%%%%%%%%%%%%%%%%%%%%%

\section{Important Signal Relationships}

\[
\boxed{
\frac{d}{dt}u(t)=\delta(t)
}
\]

\[
\boxed{
\frac{d}{dt}r(t)=u(t)
}
\]

\[
\boxed{
\frac{d^2}{dt^2}r(t)=\delta(t)
}
\]

\[
\boxed{
u(t)=\int_{-\infty}^{t}\delta(\tau)d\tau
}
\]

%%%%%%%%%%%%%%%%%%%%%%%%%%%%%%%%%%%%%%%%%%%%%%%%%%%%%%%%%%%%%%%%%%%%%%%%%%%%%%%%

\section{Quick Classification Examples}

\begin{table}[H]
\centering
\caption{Signal classification examples}
\begin{tabular}{|p{5cm}|p{6cm}|}
\hline
\textbf{Signal} & \textbf{Type} \\
\hline
\(e^{-t}u(t)\) & Energy signal \\
\hline
\(\sin(10t)\) & Power signal \\
\hline
\(u(t)\) & Power signal \\
\hline
\(\delta(t)\) & Energy signal \\
\hline
\(t^2\) & Neither energy nor power \\
\hline
\end{tabular}
\end{table}

%%%%%%%%%%%%%%%%%%%%%%%%%%%%%%%%%%%%%%%%%%%%%%%%%%%%%%%%%%%%%%%%%%%%%%%%%%%%%%%%

\section{Chapter Summary}

In this chapter, we introduced the fundamental concept of signals and classified
them into continuous-time, discrete-time, analog, digital, deterministic,
random, periodic, aperiodic, even and odd signals.

We also studied the elementary signals that form the foundation of the entire
Signals and Systems course:

\begin{itemize}
\item unit impulse,
\item unit step,
\item rectangular pulse,
\item ramp,
\item exponential,
\item sinusoidal,
\item and complex exponential signals.
\end{itemize}

Finally, we introduced the concepts of energy and power signals, which are
essential for Fourier analysis and communication theory.

%%%%%%%%%%%%%%%%%%%%%%%%%%%%%%%%%%%%%%%%%%%%%%%%%%%%%%%%%%%%%%%%%%%%%%%%%%%%%%%%

\begin{tcolorbox}[title=One-Minute Revision,colback=blue!5]

\[
x(t)\rightarrow \text{CT signal}
\]

\[
x[n]\rightarrow \text{DT signal}
\]

\[
E=\int_{-\infty}^{\infty}|x(t)|^2dt
\]

\[
P=\lim_{T\to\infty}\frac{1}{2T}\int_{-T}^{T}|x(t)|^2dt
\]

\[
\frac{d}{dt}u(t)=\delta(t)
\]

\[
r(t)=tu(t)
\]

\[
x(t)=Ae^{at}
\]

\[
x(t)=A\cos(\omega t+\phi)
\]

\[
T=\frac{2\pi}{\omega}
\]

\[
e^{j\theta}=\cos\theta+j\sin\theta
\]

\end{tcolorbox}
%%%%%%%%%%%%%%%%%%%%%%%%%%%%%%%%%%%%%%%%%%%%%%%%%%%%%%%%%%%%%%%%%%%%%%%%%%%%%%%%